\chapter{Meta-Instructions}
\label{ch:meta-instructions}

\section{Overview}

Meta-instructions are assembly mnemonics that lower to ordinary SIRCIS instruction encodings. They exist to make common
operations easier to read and write; they do not introduce additional CPU opcodes.

\begin{table}[H]
	\centering
	\begin{tabular}{lll}
		\toprule
		\textbf{Meta-Instruction} & \textbf{Primary Documentation}     & \textbf{Lowers To}          \\
		\midrule
		\mnemonic{NOOP}           & Chapter~\ref{ch:meta-instructions} & \mnemonic{ADDI[N] r1, \#0}  \\
		\mnemonic{SHFT}           & Chapter~\ref{ch:alu-instructions}  & \mnemonic{ORRI[S]}          \\
		\mnemonic{BRAN}           & Chapter~\ref{ch:control-flow}      & \mnemonic{LDEA p, ...}      \\
		\mnemonic{BRSR}           & Chapter~\ref{ch:control-flow}      & \mnemonic{LDEL p, ...}      \\
		\mnemonic{LJMP}           & Chapter~\ref{ch:control-flow}      & \mnemonic{LDEA p, ...}      \\
		\mnemonic{LJSR}           & Chapter~\ref{ch:control-flow}      & \mnemonic{LDEL p, ...}      \\
		\mnemonic{RETS}           & Chapter~\ref{ch:control-flow}      & \mnemonic{LDEA p, (\#0, l)} \\
		\mnemonic{EXCP}           & Chapter~\ref{ch:coprocessor}       & \mnemonic{COPI}             \\
		\mnemonic{WAIT}           & Chapter~\ref{ch:coprocessor}       & \mnemonic{COPI \#0x1900}    \\
		\mnemonic{RETE}           & Chapter~\ref{ch:coprocessor}       & \mnemonic{COPI \#0x1A00}    \\
		\mnemonic{RSET}           & Chapter~\ref{ch:coprocessor}       & \mnemonic{COPI \#0x1B00}    \\
		\mnemonic{ETFR}           & Chapter~\ref{ch:coprocessor}       & \mnemonic{COPI \#0x1Cxy}    \\
		\mnemonic{ETTR}           & Chapter~\ref{ch:coprocessor}       & \mnemonic{COPI \#0x1Dxy}    \\
		\mnemonic{DMAR}           & Chapter~\ref{ch:coprocessor}       & \mnemonic{COPI \#0x28xx}    \\
		\mnemonic{DMAW}           & Chapter~\ref{ch:coprocessor}       & \mnemonic{COPI \#0x29xx}    \\
		\mnemonic{DMAT}           & Chapter~\ref{ch:coprocessor}       & \mnemonic{COPI \#0x2Ann}    \\
		\mnemonic{MULU}           & Chapter~\ref{ch:coprocessor}       & \mnemonic{COPI \#0x3000}    \\
		\mnemonic{MULS}           & Chapter~\ref{ch:coprocessor}       & \mnemonic{COPI \#0x3100}    \\
		\mnemonic{DIVU}           & Chapter~\ref{ch:coprocessor}       & \mnemonic{COPI \#0x3200}    \\
		\mnemonic{DIVS}           & Chapter~\ref{ch:coprocessor}       & \mnemonic{COPI \#0x3300}    \\
		\bottomrule
	\end{tabular}
	\caption{Meta-Instruction Cross-Reference}
	\label{tab:meta-instruction-cross-reference}
\end{table}

\begin{instructionbox}{NOOP}{No Operation}

	\textbf{Applicability:} SIRC-1, all revisions

	\textbf{Assembles to:} \texttt{ADDI[N] r1, \#0}

	\textbf{Instruction Fields:}
	\begin{itemize}
		\item \textbf{Register (bits 25--22):} fixed at \texttt{0x0} (\reg{r1}).
		\item \textbf{Immediate (bits 21--6):} fixed at \imm{0}.
		\item \textbf{AF (bits 5--4):} fixed at \texttt{00} (status register not updated).
		\item \textbf{Condition field (bits 3--0):} any of the 16 condition-code encodings in
		      Table~\ref{tab:condition-code-encoding}.
	\end{itemize}

	\textbf{Syntax:} \texttt{NOOP}

	\textbf{Operands:} None.

	\textbf{Operation:} Executes an add-immediate-zero instruction whose destination is \reg{r1}: \reg{r1} is rewritten
	with its own value plus zero, so no architectural state changes occur. Because the status-register update source is
	\texttt{None}, no condition flags are written. Register field \texttt{0x0} is deliberately assigned to \reg{r1}, an
	ordinary unprivileged register (see Chapter~\ref{ch:registers}), so this instruction is unprivileged and an all-zero
	instruction word (as read from erased or uninitialized memory) always executes safely as a no-op, in both protected
	and supervisor mode.

	\textbf{Description:} Does nothing for 6 cycles. Useful for timing delays, instruction padding, or placeholders during
	development.

	\textbf{Write-back:} No architectural state changes.

	\textbf{Flags:} Preserved.

	\textbf{Exceptions:} Inherits \mnemonic{ADDI[N]} exception behavior. It does not access data memory and raises no
	instruction-specific faults; instruction-fetch faults can still occur before it executes, and trace faults can still
	occur after it commits.

	\textbf{Condition field:} If the condition is false, no state changes occur. Conditional \mnemonic{NOOP} is allowed but
	normally not useful.

	\textbf{Timing:} Same as \mnemonic{ADDI}: 6 cycles, plus instruction-fetch wait states.

	\textbf{Privilege:} Available in protected mode and supervisor mode.

	\textbf{Example 17-1:}
	\begin{lstlisting}
NOOP                          ; Wait 6 cycles
NOOP
NOOP                          ; Total 18 cycles delay
\end{lstlisting}

\end{instructionbox}
